% Imporditud: tools/import_eksperimendid.py
% Allikas: Eesti füüsikaolümpiaadi eksperimendi ülesannete
%   kogu (Rael Kalda, 2023), ülesanne Ü85, lahendus L85.
\setAuthor{EFO žürii}
\setRound{lõppvoor}
\setYear{2000}
\setNumber{G E1}
\setDifficulty{3}
\setTopic{geomeetriline optika}
\setVahendid{valge paber, pliiats, kaks tasapeeglit}
\setPunktid{12}
\setAllikas{Eksperimendi kogu Ü85, lahendus lk 70}
\setLahendusseis{toores}

\prob{Nurga joonistamine}
Joonistage valgele paberile nurk 36$^\circ$.

\hint

\solu
Kui kaks tasapeeglit on teineteise suhtes mingi nurga all peegelpinnad vastakuti ja

peeglite vahele paigutada ese, tekib peeglites esemest mitu kujutist. Kujutiste arv

sõltub peeglitevahelisest nurgast. Teatud nurkade puhul kattuvad kaks kujutist ja sel juhul saab kujutiste arvu N peeglites arvutada seosest N = 360$^\circ$/$\alpha -$ 1, kus $\alpha$ on peeglite vaheline nurk kraadides. Näiteks, kui peeglid on 90$^\circ$ nurga all, tekib kolm kujutist, 45$^\circ$ nurga puhul seitse kujutist jne.

Hindamine: Idee --- 4 p, kujutiste arv 90$^\circ$ nurga korral --- 2 p, kujutiste arv näiteks 45$^\circ$ nurga korral --- 2 p, algoritmi väljatöötamine --- 2 p, täpse nurga joonestamine

--- 2 p.
\probend
